% ============================================================
\chapter{$\sigma$-Algebras and Measurable Spaces}
\label{ch:sigma_algebras}
% ============================================================

At the end of the nineteenth century, mathematicians faced a fundamental
problem: how to rigorously define the ``size'' of a set. The Riemann
integral, which had served so well for decades, was showing its limits in
the face of pathological functions discovered by Dirichlet, Weierstrass, and
others. Henri Lebesgue, in his 1902 thesis, proposed a revolution: instead
of slicing the $x$-axis into small intervals, he sliced the $y$-axis. But
for this idea to work, one first had to answer a preliminary question: to
which sets can we assign a measure in a consistent way?

The answer --- and this is the subject of this first chapter --- is that we
must restrict ourselves to a collection of sets with good closure
properties: a $\sigma$-algebra. This concept, which may appear purely
technical at first sight, is in reality the keystone of all of measure
theory, integration, and by extension modern probability.

\section{Introduction}

The first step in building measure theory is to identify the sets to which we wish
to assign a ``size.'' In general, one cannot consistently measure \emph{all} subsets
of a given set (we shall see in Chapter~\ref{ch:lebesgue} that non-measurable sets
exist). One must therefore restrict attention to a collection of sets that is rich
enough to be useful, yet structured enough for the measure to be well-defined. This
collection is a $\sigma$-algebra.

\section{Fundamental definitions}

\begin{definition}[$\sigma$-algebra]
Let $X$ be a non-empty set. A \emph{$\sigma$-algebra} on $X$ is a collection
$\mathcal{A} \subset \mathcal{P}(X)$ of subsets of $X$ satisfying:
\begin{enumerate}
    \item[(i)] $X \in \mathcal{A}$;
    \item[(ii)] if $A \in \mathcal{A}$, then $A^c = X \setminus A \in \mathcal{A}$;
    \item[(iii)] if $(A_n)_{n \in \N}$ is a sequence of elements of $\mathcal{A}$,
    then $\bigcup_{n=0}^{\infty} A_n \in \mathcal{A}$.
\end{enumerate}
The pair $(X, \mathcal{A})$ is called a \emph{measurable space}.
\end{definition}

\begin{remark}
Condition (iii) concerns \emph{countable} unions. This is the essential difference
with the notion of an \emph{algebra} (or Boolean algebra), which only requires
closure under \emph{finite} unions.
\end{remark}

\begin{proposition}[Elementary properties]
\label{prop:sigma_props}
Let $\mathcal{A}$ be a $\sigma$-algebra on $X$. Then:
\begin{enumerate}
    \item $\varnothing \in \mathcal{A}$;
    \item $\mathcal{A}$ is closed under countable intersections: if
    $(A_n)_{n \in \N} \subset \mathcal{A}$, then
    $\bigcap_{n=0}^{\infty} A_n \in \mathcal{A}$;
    \item $\mathcal{A}$ is closed under set difference: if $A, B \in \mathcal{A}$,
    then $A \setminus B \in \mathcal{A}$;
    \item $\mathcal{A}$ is closed under finite unions and intersections.
\end{enumerate}
\end{proposition}

\begin{proof}
(1) $\varnothing = X^c \in \mathcal{A}$ by (i) and (ii).

(2) By De Morgan's laws:
$\bigcap_n A_n = \bigl(\bigcup_n A_n^c\bigr)^c$. Since each
$A_n^c \in \mathcal{A}$ by (ii), the union $\bigcup_n A_n^c \in \mathcal{A}$
by (iii), and its complement is in $\mathcal{A}$ by (ii).

(3) $A \setminus B = A \cap B^c$, which is in $\mathcal{A}$ by (ii) and (2).

(4) Special case of (iii) and (2) with $A_n = \varnothing$ for $n$ large enough.
\end{proof}

\begin{definition}[Algebra]
A collection $\mathcal{A}_0 \subset \mathcal{P}(X)$ is an \emph{algebra} on $X$
if it satisfies (i), (ii), and closure under \emph{finite} unions:
if $A, B \in \mathcal{A}_0$, then $A \cup B \in \mathcal{A}_0$.
\end{definition}

\begin{example}
Here are some fundamental examples of $\sigma$-algebras:
\begin{enumerate}
    \item The \emph{trivial $\sigma$-algebra} $\{\varnothing, X\}$ is the smallest
    $\sigma$-algebra on $X$.
    \item The \emph{discrete $\sigma$-algebra} $\mathcal{P}(X)$ is the largest
    $\sigma$-algebra on $X$.
    \item For any $A \subset X$, $\mathcal{A} = \{\varnothing, A, A^c, X\}$ is a
    $\sigma$-algebra.
    \item If $X$ is uncountable, the collection of countable sets and sets with
    countable complement forms a $\sigma$-algebra, called the
    \emph{co-countable $\sigma$-algebra}.
\end{enumerate}
\end{example}

\section{Generated $\sigma$-algebra}

In practice, one does not define a $\sigma$-algebra by listing all its elements, but
by ``generating'' it from a smaller collection.

\begin{theorem}[Generated $\sigma$-algebra]
\label{thm:generated_sigma}
Let $\mathcal{E} \subset \mathcal{P}(X)$ be any collection of subsets of $X$. There
exists a smallest $\sigma$-algebra containing $\mathcal{E}$, called the
\emph{$\sigma$-algebra generated by $\mathcal{E}$} and denoted $\sigma(\mathcal{E})$.
It is given by:
\[
\sigma(\mathcal{E}) = \bigcap_{\substack{\mathcal{A} \text{ is a $\sigma$-algebra} \\
\mathcal{E} \subset \mathcal{A}}} \mathcal{A}.
\]
\end{theorem}

\begin{proof}
The intersection is well-defined since at least $\mathcal{P}(X)$ is a
$\sigma$-algebra containing $\mathcal{E}$, so the family is non-empty.

We verify that $\sigma(\mathcal{E})$ is a $\sigma$-algebra:
\begin{itemize}
    \item $X$ belongs to every $\sigma$-algebra, so $X \in \sigma(\mathcal{E})$.
    \item If $A \in \sigma(\mathcal{E})$, then $A$ belongs to every $\sigma$-algebra
    $\mathcal{A} \supset \mathcal{E}$, so $A^c$ does too, hence
    $A^c \in \sigma(\mathcal{E})$.
    \item The same argument applies to countable unions.
\end{itemize}
Finally, $\sigma(\mathcal{E})$ is the \emph{smallest} such $\sigma$-algebra since
it is contained in every $\sigma$-algebra containing $\mathcal{E}$ (by definition
as an intersection).
\end{proof}

\begin{remark}
The intersection construction is non-constructive: it does not allow one to describe
the elements of $\sigma(\mathcal{E})$ explicitly. In general,
$\sigma(\mathcal{E})$ is much larger than $\mathcal{E}$ and its structure is complex.
\end{remark}

\section{The Borel $\sigma$-algebra}

\begin{definition}[Borel $\sigma$-algebra]
Let $(X, \tau)$ be a topological space. The \emph{Borel $\sigma$-algebra} of $X$,
denoted $\mathcal{B}(X)$, is the $\sigma$-algebra generated by the open sets of $X$:
\[
\mathcal{B}(X) = \sigma(\tau).
\]
The elements of $\mathcal{B}(X)$ are called \emph{Borel sets}.
\end{definition}

\begin{theorem}[Generators of the Borel $\sigma$-algebra of $\R$]
\label{thm:borel_generators}
The Borel $\sigma$-algebra $\mathcal{B}(\R)$ is generated by each of the following
collections:
\begin{enumerate}
    \item the open subsets of $\R$;
    \item the closed subsets of $\R$;
    \item the open intervals $(a, b)$, $a < b$;
    \item the half-open intervals $[a, b)$, $a < b$;
    \item the half-open intervals $(a, b]$, $a < b$;
    \item the rays $(-\infty, a)$, $a \in \R$;
    \item the rays $(-\infty, a]$, $a \in \R$;
    \item the rays $(a, +\infty)$, $a \in \R$.
\end{enumerate}
\end{theorem}

\begin{proof}
Denote by $\mathcal{B}$ the Borel $\sigma$-algebra and by $\sigma_i$ the
$\sigma$-algebra generated by the $i$-th collection. It suffices to show
circular inclusions.

$\sigma_6 \subset \sigma_7$: $(-\infty, a) = \bigcup_{n=1}^{\infty}
(-\infty, a - 1/n]$, so $(-\infty, a) \in \sigma_7$.

$\sigma_7 \subset \sigma_6$: $(-\infty, a] = \bigcap_{n=1}^{\infty}
(-\infty, a + 1/n)$, so $(-\infty, a] \in \sigma_6$.

$\sigma_3 \subset \sigma_6$: $(a, b) = (-\infty, b) \cap (a, +\infty)$. Now
$(a, +\infty) = (-\infty, a]^c \in \sigma_6$ (since $\sigma_6 = \sigma_7$ as shown).

$\sigma_1 \subset \sigma_3$: every open subset of $\R$ is a countable union of
open intervals (Lindel\"of property).

Since $\sigma_1 = \mathcal{B}$ by definition and
$\sigma_1 \subset \sigma_3 \subset \sigma_6 \subset \mathcal{B}$, we have equality.
The remaining cases are analogous.
\end{proof}

\begin{keyformulas}[title=\textbf{Common generators of $\mathcal{B}(\R)$}]
\[
\mathcal{B}(\R) = \sigma\bigl(\{(a,b) : a < b\}\bigr)
= \sigma\bigl(\{(-\infty, a) : a \in \R\}\bigr)
= \sigma\bigl(\{(-\infty, a] : a \in \R\}\bigr).
\]
\end{keyformulas}

\section{Borel $\sigma$-algebra on $\R^n$}

\begin{definition}
The Borel $\sigma$-algebra on $\R^n$ is $\mathcal{B}(\R^n) = \sigma(\tau_n)$, where
$\tau_n$ is the standard topology on $\R^n$.
\end{definition}

\begin{proposition}
We have $\mathcal{B}(\R^n) = \sigma(\mathcal{R}_n)$, where $\mathcal{R}_n$ is the
collection of open rectangles
$\prod_{i=1}^n (a_i, b_i)$ with $a_i < b_i$ for all $i$.
\end{proposition}

\begin{proof}
Every open set in $\R^n$ is a countable union of open rectangles with rational
vertices (since $\Q^{2n}$ is countable), so
$\sigma(\tau_n) \subset \sigma(\mathcal{R}_n)$. Conversely, each open rectangle is
open in $\R^n$, so $\sigma(\mathcal{R}_n) \subset \sigma(\tau_n)$.
\end{proof}

\section{Product $\sigma$-algebra}

\begin{definition}[Product $\sigma$-algebra]
Let $(X_1, \mathcal{A}_1)$ and $(X_2, \mathcal{A}_2)$ be measurable spaces. The
\emph{product $\sigma$-algebra} is:
\[
\mathcal{A}_1 \otimes \mathcal{A}_2 = \sigma\bigl(\{A_1 \times A_2 :
A_1 \in \mathcal{A}_1,\, A_2 \in \mathcal{A}_2\}\bigr).
\]
More generally, for a family $(X_i, \mathcal{A}_i)_{i \in I}$:
\[
\bigotimes_{i \in I} \mathcal{A}_i = \sigma\bigl(\{\pi_i^{-1}(A_i) :
i \in I,\, A_i \in \mathcal{A}_i\}\bigr),
\]
where $\pi_i : \prod_{j \in I} X_j \to X_i$ is the canonical projection.
\end{definition}

\begin{theorem}
\label{thm:borel_product}
If $X_1$ and $X_2$ are separable metric spaces, then
\[
\mathcal{B}(X_1) \otimes \mathcal{B}(X_2) = \mathcal{B}(X_1 \times X_2).
\]
In particular, $\mathcal{B}(\R^n) = \mathcal{B}(\R) \otimes \cdots \otimes
\mathcal{B}(\R)$ ($n$ times).
\end{theorem}

\begin{proof}
Let $U$ be an open set in $X_1 \times X_2$. For every $(x_1, x_2) \in U$, there
exist open sets $V_1 \ni x_1$ and $V_2 \ni x_2$ with $V_1 \times V_2 \subset U$.
By separability, one can choose $V_1$ and $V_2$ from countable bases
$\mathcal{B}_1$ and $\mathcal{B}_2$ of $X_1$ and $X_2$. Thus
$U = \bigcup_k V_1^{(k)} \times V_2^{(k)}$, a countable union of measurable
rectangles. This shows $\mathcal{B}(X_1 \times X_2) \subset
\mathcal{B}(X_1) \otimes \mathcal{B}(X_2)$.

Conversely, if $A_1 \in \mathcal{B}(X_1)$, then $\pi_1^{-1}(A_1) =
A_1 \times X_2$ is Borel in $X_1 \times X_2$ since $\pi_1$ is continuous. Similarly
for $\pi_2$. The sets $A_1 \times A_2 = \pi_1^{-1}(A_1) \cap \pi_2^{-1}(A_2)$ are
Borel, giving the reverse inclusion.
\end{proof}

\section{Sub-$\sigma$-algebra and trace}

\begin{definition}[Trace]
Let $(X, \mathcal{A})$ be a measurable space and $Y \subset X$. The \emph{trace} of
$\mathcal{A}$ on $Y$ is:
\[
\mathcal{A}|_Y = \{A \cap Y : A \in \mathcal{A}\}.
\]
This is a $\sigma$-algebra on $Y$.
\end{definition}

\begin{proposition}
If $\mathcal{A} = \sigma(\mathcal{E})$, then
$\mathcal{A}|_Y = \sigma(\mathcal{E}|_Y)$, where
$\mathcal{E}|_Y = \{E \cap Y : E \in \mathcal{E}\}$.
\end{proposition}

\begin{proof}
Set $\mathcal{F} = \{A \subset X : A \cap Y \in \sigma(\mathcal{E}|_Y)\}$.
One checks that $\mathcal{F}$ is a $\sigma$-algebra containing $\mathcal{E}$, hence
$\sigma(\mathcal{E}) \subset \mathcal{F}$, giving
$\mathcal{A}|_Y \subset \sigma(\mathcal{E}|_Y)$. The reverse inclusion is immediate.
\end{proof}

\section{Measurable maps}

\begin{definition}[Measurable map]
Let $(X, \mathcal{A})$ and $(Y, \mathcal{B})$ be measurable spaces. A map
$f : X \to Y$ is \emph{$(\mathcal{A}, \mathcal{B})$-measurable} (or simply
\emph{measurable}) if:
\[
\forall\, B \in \mathcal{B}, \quad f^{-1}(B) \in \mathcal{A}.
\]
\end{definition}

\begin{proposition}[Generator criterion for measurability]
\label{prop:measurability_criterion}
Let $\mathcal{B} = \sigma(\mathcal{E})$. Then $f : (X, \mathcal{A}) \to
(Y, \mathcal{B})$ is measurable if and only if:
\[
\forall\, E \in \mathcal{E}, \quad f^{-1}(E) \in \mathcal{A}.
\]
\end{proposition}

\begin{proof}
The forward direction is clear. For the converse, set
$\mathcal{C} = \{B \subset Y : f^{-1}(B) \in \mathcal{A}\}$.
We verify that $\mathcal{C}$ is a $\sigma$-algebra:
\begin{itemize}
    \item $f^{-1}(Y) = X \in \mathcal{A}$, so $Y \in \mathcal{C}$.
    \item If $B \in \mathcal{C}$, then $f^{-1}(B^c) = f^{-1}(B)^c \in \mathcal{A}$,
    so $B^c \in \mathcal{C}$.
    \item If $(B_n) \subset \mathcal{C}$, then
    $f^{-1}(\bigcup_n B_n) = \bigcup_n f^{-1}(B_n) \in \mathcal{A}$, so
    $\bigcup_n B_n \in \mathcal{C}$.
\end{itemize}
By hypothesis, $\mathcal{E} \subset \mathcal{C}$, hence
$\mathcal{B} = \sigma(\mathcal{E}) \subset \mathcal{C}$, meaning $f$ is measurable.
\end{proof}

\begin{corollary}
A map $f : \R \to \R$ is Borel measurable (i.e.\
$(\mathcal{B}(\R), \mathcal{B}(\R))$-measurable) if and only if
$f^{-1}((-\infty, a)) \in \mathcal{B}(\R)$ for every $a \in \R$.
\end{corollary}

\begin{proposition}[Composition]
If $f : (X, \mathcal{A}) \to (Y, \mathcal{B})$ and
$g : (Y, \mathcal{B}) \to (Z, \mathcal{C})$ are measurable, then
$g \circ f : (X, \mathcal{A}) \to (Z, \mathcal{C})$ is measurable.
\end{proposition}

\begin{proof}
For any $C \in \mathcal{C}$, $(g \circ f)^{-1}(C) = f^{-1}(g^{-1}(C))$. Since $g$
is measurable, $g^{-1}(C) \in \mathcal{B}$, and since $f$ is measurable,
$f^{-1}(g^{-1}(C)) \in \mathcal{A}$.
\end{proof}

\section{Dynkin systems and the $\pi$-$\lambda$ lemma}

Dynkin's $\pi$-$\lambda$ lemma is a fundamental tool for proving that two measures
agree.

\begin{definition}[$\pi$-system and $\lambda$-system]
\begin{itemize}
    \item A \emph{$\pi$-system} on $X$ is a collection $\mathcal{P}
    \subset \mathcal{P}(X)$ closed under finite intersections.
    \item A \emph{$\lambda$-system} (or \emph{Dynkin system}) on $X$ is a collection
    $\mathcal{D} \subset \mathcal{P}(X)$ such that:
    \begin{enumerate}
        \item[(i)] $X \in \mathcal{D}$;
        \item[(ii)] if $A, B \in \mathcal{D}$ and $A \subset B$, then
        $B \setminus A \in \mathcal{D}$;
        \item[(iii)] if $(A_n)$ is an increasing sequence in $\mathcal{D}$, then
        $\bigcup_n A_n \in \mathcal{D}$.
    \end{enumerate}
\end{itemize}
\end{definition}

\begin{theorem}[Dynkin's $\pi$-$\lambda$ lemma]
\label{thm:pi_lambda}
If $\mathcal{P}$ is a $\pi$-system and $\mathcal{D}$ is a $\lambda$-system with
$\mathcal{P} \subset \mathcal{D}$, then $\sigma(\mathcal{P}) \subset \mathcal{D}$.
\end{theorem}

\begin{proof}
It suffices to show that $d(\mathcal{P})$, the smallest $\lambda$-system containing
$\mathcal{P}$, is a $\sigma$-algebra (for then
$\sigma(\mathcal{P}) \subset d(\mathcal{P}) \subset \mathcal{D}$).

For this, it suffices to show that $d(\mathcal{P})$ is closed under finite
intersections (since a $\lambda$-system closed under finite intersections is a
$\sigma$-algebra).

For $A \in d(\mathcal{P})$, set
$\mathcal{G}_A = \{B \in d(\mathcal{P}) : A \cap B \in d(\mathcal{P})\}$.

\emph{Step 1.} If $A \in \mathcal{P}$, then $\mathcal{G}_A$ is a $\lambda$-system
containing $\mathcal{P}$ (since $\mathcal{P}$ is a $\pi$-system), hence
$\mathcal{G}_A \supset d(\mathcal{P})$.

\emph{Step 2.} For any $A \in d(\mathcal{P})$, $\mathcal{G}_A$ is a
$\lambda$-system. By Step 1, $\mathcal{P} \subset \mathcal{G}_A$ for every
$A \in d(\mathcal{P})$, so $\mathcal{G}_A \supset d(\mathcal{P})$.

Thus $d(\mathcal{P})$ is closed under finite intersections, which concludes.
\end{proof}

\begin{warning}[title=\textbf{Using the $\pi$-$\lambda$ lemma}]
This lemma is very useful for showing that a property true on a $\pi$-system extends
to the generated $\sigma$-algebra. Typically, to show that two measures $\mu$ and
$\nu$ on $\sigma(\mathcal{P})$ agree, it suffices to verify $\mu = \nu$ on the
$\pi$-system $\mathcal{P}$ (under a $\sigma$-finiteness condition).
\end{warning}

\section{Uniqueness of measures}

\begin{theorem}[Uniqueness of measures]
\label{thm:uniqueness_measures}
Let $\mu$ and $\nu$ be two measures on $(X, \sigma(\mathcal{P}))$, where
$\mathcal{P}$ is a $\pi$-system. Suppose there exists a sequence
$(E_n)_{n \geq 1} \subset \mathcal{P}$ with $E_n \uparrow X$ and
$\mu(E_n) = \nu(E_n) < \infty$ for all $n$. If $\mu = \nu$ on $\mathcal{P}$,
then $\mu = \nu$ on $\sigma(\mathcal{P})$.
\end{theorem}

\begin{proof}
Fix $n$ and set $\mu_n(\cdot) = \mu(\cdot \cap E_n)$ and
$\nu_n(\cdot) = \nu(\cdot \cap E_n)$. These are finite measures. The collection
\[
\mathcal{D}_n = \{A \in \sigma(\mathcal{P}) : \mu_n(A) = \nu_n(A)\}
\]
is a $\lambda$-system (direct verification). By hypothesis,
$\mathcal{P} \subset \mathcal{D}_n$. By the $\pi$-$\lambda$ lemma,
$\sigma(\mathcal{P}) \subset \mathcal{D}_n$, so $\mu_n = \nu_n$ on
$\sigma(\mathcal{P})$.

For every $A \in \sigma(\mathcal{P})$, $A \cap E_n \uparrow A$, so by monotone
continuity:
$\mu(A) = \lim_n \mu_n(A) = \lim_n \nu_n(A) = \nu(A)$.
\end{proof}

\section{Exercises}

\begin{exercise}
Let $X$ be an uncountable set. Show that the collection
\[
\mathcal{A} = \{A \subset X : A \text{ is countable or } A^c \text{ is countable}\}
\]
is a $\sigma$-algebra on $X$.
\end{exercise}

\begin{exercise}
Show that $\mathcal{B}(\R)$ is generated by the open intervals with rational
endpoints.
\end{exercise}

\begin{exercise}
Let $f : \R \to \R$ be a continuous function. Show that $f$ is Borel measurable.
\end{exercise}

\begin{exercise}
Show that $\mathcal{B}(\R) \neq \mathcal{P}(\R)$.
\emph{Hint}: use a cardinality argument ($\abs{\mathcal{B}(\R)} = \abs{\R} =
\mathfrak{c}$ while $\abs{\mathcal{P}(\R)} = 2^{\mathfrak{c}}$).
\end{exercise}

\begin{exercise}
Let $(X, \mathcal{A})$ be a measurable space and $f, g : X \to \R$ two measurable
functions. Show that $\{x \in X : f(x) < g(x)\} \in \mathcal{A}$.
\end{exercise}

\begin{exercise}[Atoms of a $\sigma$-algebra]
Let $(X, \mathcal{A})$ be a measurable space. For $x \in X$, define the
\emph{atom} of $x$ as $[x] = \bigcap_{A \in \mathcal{A},\, x \in A} A$.
\begin{enumerate}
    \item Show that if $\mathcal{A}$ is generated by a countable partition, then
    $[x] \in \mathcal{A}$ for every $x$.
    \item Give an example where $[x] \notin \mathcal{A}$.
\end{enumerate}
\end{exercise}
